\begin{abstract}
% Context and problem: supported observations, not the author's conclusion.
Cognitive behavioral therapy (CBT) has evidence of short-term benefit for adult attention-deficit/hyperactivity disorder (ADHD), but the certainty and interpretation of that evidence depend on comparator and study design \citep{Lopez2018}.
The motivating question here is whether baseline executive function identifies who benefits, rather than whether executive function can improve during treatment.
Those are different estimands: prediction concerns later outcome, whereas moderation concerns variation in a treatment contrast \citep{Kraemer2002}.
% Objective and approach: prospective work, not completed results.
We propose a meta-review of systematic reviews, supplemented by original adult-ADHD studies measuring baseline executive function and subsequent CBT outcomes.
The protocol separates cognitive restructuring from multicomponent packages, prognostic associations from treatment interactions, and active from passive comparators.
It specifies ten search slices across Europe PMC, PubMed and PsycINFO, independent screening, review-overlap mapping and structured extraction of estimates and uncertainty.
% Current deliverable and limit.
This initial report provides an executed Europe PMC count baseline, an auditable citation ledger and preserved author drafts; it does not report completed screening or a pooled treatment effect.
The executive-demand critique remains a hypothesis requiring direct tests, alongside evidence for benefit, null findings and alternative explanations.
\end{abstract}
